\chapter{About probability}
\label{ch:prob}

\index{probability}

\textbf{Probability theory} is an area of mathematics that deals with
\textit{quantifying uncertainty}. We may not be sure that something is going to
happen, but we want to be precise about \textit{how} sure we are, given
everything else we may know.

\index{random value}
\index{distribution}

It turns out that probability theory is a crucial concept in Data Science.
We've already used some of it in Chapters~\ref{ch:random} and
\ref{ch:synthetic} when we generated random values from distributions. In this
chapter, we'll zoom out from programming language details and consider how to
reason about the probabilities of various values occurring and the
relationships between random variables.

\section{Terms}

\index{outcome (probability)}
\index{redheads}
\index{Girl on the Train, The@\textit{Girl on the Train, The}}
\index{Hawkins, Paula}
\index{hardback edition}

Some terms: an \textbf{outcome} is something that may (or may not) happen. This
varies widely depending on the domain: a child that's born may or may not be
male; an airline passenger may or may not have red hair; a customer may or may
not buy a particular product. ``Male,'' ``red,'' and ``\textit{The Girl on the
Train}, hardback edition'' are all outcomes in these examples.

\index{event (probability)}
\index{DVD}

An \textbf{event}, by contrast, is a \textit{set} (group) of outcomes that
we're interested in. Often we don't care so much whether or not the customer
specifically bought ``\textit{The Girl on the Train}, hardback edition'' as
whether they bought \textit{any} book by Paula Hawkins, or maybe any hardback
book, or maybe any book at all (as opposed to a DVD).

\index{sample space (probability)}
\index{redheads}
\index{bald}
\index{brunettes}
\index{blondes}
\index{omega ($\Omega$)}
\index{$\Omega$ (omega)}

Sometimes we'll use the symbol $\Omega$ (``omega'') to refer to the
\textbf{sample space}, which is simply the set of \textit{all possible
outcomes}. This will vary depending on what's under consideration; it might be
\{ male, female \} in one case and \{ red, blonde, brunette, bald, other \} in
another case.

We'll also use the notation $\overline{\phantom{xx}}$ (overbar) to mean ``any
outcome \textit{except} something.'' For instance, $\overline{\textrm{red}}$
means ``the airline passenger had a hair color \textit{other than} red.''

\section{Probability measures}

\index{probability measure}

A \textbf{probability measure} is a function that tells us intuitively how
likely an outcome is to occur. We'll use the simple notation ``P($\cdot$)'' for
this\footnote{Some authors use ``Pr()'' instead of ``P()'' for this. Same
diff.}, read as ``the probability of (something).'' The probability of any
outcome is always \textit{a number between 0 and 1}. 0 means it can never
possibly occur, 1 means it will definitely occur, and there's a continuum in
between. In our first example, we might say P(male) = .5, P(red) = .1, and
P(\textit{The Girl on the Train}, hardback edition) = .0001.

You can talk about the probability of an event too, of course, and it's simply
the sum of the outcome probabilities. P(male-or-female) = 1, P(red or blonde
or bald) = .6, P(any hardcover book) = .08. We'll say that P($\Omega$) will be
1: after all, \textit{something} has to happen, even if it's ``the customer
didn't buy anything,'' an outcome to which we'll assign a probability.

\section{``Joint'' probability}
\index{redheads}
\index{joint probability}
\index{,@, (comma) (joint probability)}
%\index{pipe@\|\ (pipe) (conditional probability)}

Very commonly, we'll be considering the likelihood that two different events
\textit{both} occur. For instance, perhaps we're interested in the probability
that an airline passenger is a red-headed female. We'll write this as P(red,
female).\footnote{You may also run across the notation
``P(red\textbf{\^{}}female),'' which means the same thing.} The comma signals
\textbf{joint probability}, which simply means that both of these
outcomes/events must occur in order for the combined event to be considered to
have occurred. If we wrote P(red, $\overline{\textrm{female}}$), it would mean
`the probability that the passenger is a red-headed \textit{non}-female.''

If you think about it, you'll realize that the joint probability cannot be
greater than either of the two individual probabilities. Adding another
condition (not female) to the first (red-headed) can't possibly
\textit{increase} the chance of it happening. Other than that, we don't know
how to compute the joint probability from the individual ones, though (yet).

Finally, note that P(A,B) is the same thing as P(B,A), guaranteed. Switching
the order of terms in a \textit{joint} probability doesn't change the answer.
(This will \textit{not} be true in the next section, so watch out!)

\section{``Conditional'' probability}

\index{conditional probability}
\index{college degree}
\index{unemployed}

Another notion we'll use a great deal is \textbf{conditional probability},
written with a ``pipe'' (vertical bar) like P(unemployed $|$\ college degree),
or P(red $|\ \overline{\textrm{female}}$). The second of these means ``if I
know for a fact that the passenger is not female, what's the probability that
they're red-headed?'' Notice this is very different than the joint probability.
With conditional probability, we're assuming that the second event
\textit{does} occur, and then asking how that revises our original prediction.
Unlike with joint probability, adding a vertical bar \textit{can} make the
answer increase.

\index{pinterest}

Consider this example: what's the probability that a random U.S. citizen has a
Pinterest account? Let's estimate that as P(pinterest) = .15. But what if we
learned that the chosen person was female? What is P(pinterest $|$ female)?
Just to shamelessly stereotype, I'm going to go out on a limb and say it's
higher; perhaps .26. On the other hand, what if we learned the person was male?
What is P(pinterest $|$ male)? I'll say it's more like .04. In both cases, we
adjusted our estimated probability in light of additional information: up in
one case, and down in the other.

\index{Bayesian reasoning}

This adjusting-our-estimates stuff is at the heart of what's called
``\textbf{Bayesian reasoning}'' (you'll see why in section~\ref{sec:bayes},
below).

\index{prior (probability)}
\index{posterior (probability)}

Two other terms that come up frequently are ``\textbf{prior probability}'' and
``\textbf{posterior probability}''. The ``\textbf{prior}'' is the probability of
an event \textit{without} applying conditioning. In the previous example .15
was the ``prior probability'' that the person was a pinterest user. When we
learned our subject was female, we revised that to the ``\textbf{posterior}''
of .26.

\index{NBA}
\index{tall}

Finally, note that P(A$|$B) is \textbf{\textit{not}} the same thing as
P(B$|$A). I gave you a heads-up about this above. In general, P(A$|$B) and
P(B$|$A) are very different numbers. Ask yourself this: if you chose a random
American, and they turned out to be pretty tall (say, taller than average),
what's the probability that they're also an NBA player? I think you'll agree
that the number here would be very small indeed: perhaps P(NBA$|$tall) would be
.00001. There just aren't that many NBA players out there, so the chances of
you randomly choosing one -- tall or not -- is miniscule. But consider the
reverse question: if your random choice just happened to be an NBA player,
what's the probability that they'd also be tall? This, by contrast, would be
far greater: I'll bet P(tall$|$NBA) is close to .95, since nearly all
professional basketball players are pretty tall.

Or continue the pinterest example. P(pinterest $|$ female) might be .26
(\textit{i.e.}, 26\% of all females are pinterest users), but P(female $|$
pinterest) might be .87 (87\% of all pinterest users are females). We'll
quantify all this in the next section.

%\section{The relationship between joint and conditional}
%
%\index{joint probability}
%\index{conditional probability}
%
%Mathematically, here's how these two concepts (joint probability and
%conditional probability) are related. If A and B are two events, then:
%
%\vspace{-.2in}
%\begin{align*}
%\boxed{P(A|B) = \frac{P(A,B)}{P(B)}}
%\end{align*}
%
%Suppose I knew that the probability of a random U.S. citizen being a
%\textit{male pinterest user} was .04. And now suppose I knew that about half
%of U.S. citizens are male. Then, I could compute:
%
%\vspace{-.2in}
%\begin{align*}
%P(\textrm{male,pinterest}) = .04 \\
%P(\textrm{male}) = .5 \\
%P(\textrm{pinterest} | \textrm{male}) &= \frac{P(\textrm{male},
%\textrm{pinterest})}{P(\textrm{pinterest})} \\
%                    &= \frac{P(\textrm{pinterest},
%\textrm{male})}{P(\textrm{pinterest})} \\
%                    &= \frac{.04}{.15} = .267 \\
%\end{align*}
%
%\section{Bayes' Rule}
%\label{sec:bayes}
%
%Knowing how joint and conditional probabilities are related, and knowing that
%P(A,B) = P(B,A), we can derive the following law:
%
%\vspace{-.2in}
%\begin{align*}
%P(A|B) = \frac{P(B|A) \cdot P(A)}{P(B)}
%\end{align*}
%
%This is called \textbf{Bayes' Rule} (sometimes \textbf{Bayes' Law} or
%\textbf{Bayes' Theorem}) after the 18th-century Presbyterian minister Thomas
%Bayes. It is one of the most influential (and controversial, actually)
%statements in statistics and even all of mathematics.
%
%\section{Independence}
%
%We need one more background concept before we attack Na\"{i}ve Bayes Classifiers,
%and that is the \textbf{independence} of events.
%
%Simply put, two events are ``independent'' if \textit{whether one of them
%occurs has no bearing on whether the other one occurs}. It does \textit{not}
%mean the events are mutually exclusive. It means that the probability of each
%one of them occurring has nothing to do with the other.
%
%It's best seen through an example. Suppose I chose a random U.S. citizen, and
%I observe two things: whether or not they're white, and whether or not they're
%female. If I asked you for your estimate of the ``prior'' for the race
%question, you'd probably come up with something like P(white) = .6. Now I ask
%you: suppose, before observing their race, that you learned your subject was
%female. Does this cause you to update your ``prior'' estimate to a more
%accurate ``posterior''?
%
%The answer is no. Whether or not someone is white is \textit{independent} of
%whether they're female. You thought P(white) was .6, and you still think
%P(white$|$female) is .6. The opposite was true in the pinterest example: your
%estimate of their pinterest status \textit{did} change when you learned about
%their gender. Therefore, those two events are \textit{not} independent.
%
%Both of the following mathematical identities are true \textit{\textbf{only
%\underline{if}}} the events $A$ and $B$ are independent:
%\vspace{-.2in}
%\begin{align*}
%P(A,B) &= P(A) \cdot P(B) \\
%P(A|B) &= P(A)
%\end{align*}
%
